\chapter*{Preface}
\addcontentsline{toc}{chapter}{Preface}
\markboth{Preface}{Preface}

\section*{The cover: a universe of epicycles}

On the cover, three concentric circles nest quietly within one another. On the outermost ring sixty points are evenly distributed, like the graduations of a clock, hinting at the recurrence of the sexagenary cycle; one layer inwards, twelve red points mark the positions of the Earthly Branches; further inwards, eight green points correspond to the arrangement of the Eight Trigrams; and on the innermost ring, five orange points hint at the workings of the Five Phases. At the centre of the circle, a red point extends a fine line that joins the first small circle; the point on that circle in turn extends a smaller circle; and on the smaller circle there is a circle smaller still. This is a picture of epicycles---for every single step forward of the outer circle, the inner circle turns through one full revolution. The ancient Greeks used such nested circles to describe the motions of the planets, and this book will use them to describe things that possess a periodic structure \textbf{under appropriate conditions}.

At the bottom of the cover there are two lines of formulae. The first tells us that the phase point of a system with period $T$ can be labelled by the triple $(A_k, \omega_k, \phi_k)$---amplitude, frequency, phase. The second expresses the preliminary idea of periodic convolution: the deep similarity between two functions depends not only on whether their respective frequencies coincide, but also on the interaction between every pair of frequencies. \textbf{It must be stressed} that this statement is at present a \textbf{heuristic picture} rather than a mathematical theorem; the conditions required for it to hold rigorously are to be found in the \textbf{Introduction}, among the convergence, symmetry and boundedness requirements listed after the definition of periodic convolution.

This picture is a visual summary of the entire contents of the book: periodic systems can be described by families of epicycles, epicycles are characterised by amplitude, frequency and phase, and the relations between periods are described by periodic convolution. \textbf{It must be stressed once again} that this summary is a \textbf{research-programme picture}: throughout the book it is gradually converted into a number of precise, conditional mathematical propositions, rather than being a theorem obtained at a single stroke.

\section*{Periodicity: an early language for understanding the world}

Human awareness of periodicity can be traced back to the observation of the firmament. The alternation of day and night, the waxing and waning of the moon, the return of the stars---these astronomical rhythms recur with a predictable regularity, and in several early civilisations they constituted the starting point of the awareness of time. In Mesopotamia the Sumerians compiled the appearances and disappearances of Venus into star catalogues; in the Nile valley the Egyptians used the heliacal rising of Sirius to forecast the flood; on the two banks of the Yellow River the forebears measured the lengthening and shortening of the sun's shadow with the gnomon, determined the solar terms and regulated the farming seasons. Periodicity is one of the earliest tools with which human reason sought to characterise time.

The periodicities above the firmament not only gave rise to calendars, but in several civilisations also promoted the development of mathematics. When the forebears discovered that the motions of the sun and the moon are not isolated phenomena, but can be referred to one another and predicted from one another, a recognition gradually took shape: within a certain accuracy, the motions of the heavenly bodies can be grasped by rhythmic mathematical models. This recognition was one of the factors behind the independent emergence of a tradition of exact science in several civilisations---but it was not the only factor, nor did every civilisation develop this recognition in the same way.

\section*{The mathematisation of periodicity}

Once periodicity has been perceived, it is gradually mathematised. The Babylonians turned the periods of the sun, the moon and the five planets into sexagesimal numerical tables; the Maya used the nesting of multiple calendars to track the interleaving of Venus and the lunar syzygies; the ancient Greeks decomposed the motions of the heavenly bodies into composites of epicycles and deferents. In several early civilisations astronomy was one of the first branches of science to attain high precision---because the periodic motions of the heavenly bodies provided relatively stable objects of measurement.

In China this process of mathematisation unfolded in a distinctive way. The `Tianwen xun' (Astronomy) chapter of the \textit{Huainanzi}, in the words `in seventy-six years the sun and the moon return, in the first month, to the fifth degree of Yingshi, with no fractional remainder, and this is called one epoch', establishes a high-precision common-multiple approximation to the tropical year and the synodic month; the \textit{Huangji Jingshi}, by means of its nested system of yuan, hui, yun and shi---one yuan comprises twelve hui, one hui thirty yun, one yun twelve shi, and one shi thirty years---brings the periods of celestial phenomena and the periods of human affairs into one and the same self-consistent arithmetical framework; and the division into the twenty-four solar terms is an independent invention of the solar calendar: it divides one circuit of the ecliptic into twenty-four equally spaced phase nodes, each node corresponding to the approximate position of the sun on the ecliptic. \textbf{It should be pointed out} that the historical account given above belongs to traditional calendrical reckoning and to the interpretation of cultural history; its specific numerical values and its dating should be regarded as \textbf{symbolic settings or as following traditional accounts}, and should not be taken as premises for strict historical fact or for mathematical derivation.

By the seventeenth century, Kepler had turned more than twenty years of Tycho's observational data into three laws. The elliptical orbit of a planet is no longer an arbitrarily posited geometrical construction, but the necessary outcome, in an inverse-square force field, of the joint constraints of conservation of angular momentum and conservation of energy. Newton went further and unified these three laws into the equation of universal gravitation, and from then on the mathematical description of periodicity acquired the foundation of classical mechanics. \textbf{It should be pointed out} that the explanation of periodic motion given by Kepler's laws and by Newtonian mechanics holds \textbf{within the framework of classical mechanics}; within the frameworks of general relativity or of quantum mechanics the corresponding statements require modification.

\section*{From invariance to conservation laws, and from conservation laws to symmetry}

One core feature of periodicity is \textbf{invariance}. The sun rises and sets, and today resembles yesterday; the moon wanes and waxes, and this cycle resembles the one before; spring departs and autumn arrives, and this year resembles the year before. What the forebears perceived first was that certain properties remain approximately unchanged amidst change. The sun rises in roughly the same direction every day, the moon completes one cycle of waning and waxing about every twenty-nine days, and the four seasons succeed one another each year in a fixed order. These approximate invariances are the earliest knowledge of periodicity.

When invariance is rendered precise as a mathematical proposition, conservation laws make their entrance. When a quantity remains unchanged throughout a certain class of evolutionary processes, we say that this quantity is conserved in that class of processes. Kepler discovered that, as a planet moves on its elliptical orbit, the line joining the planet to the sun sweeps out equal areas in equal times---this is a special case of conservation of angular momentum. Galileo discovered that the period of a pendulum at small amplitudes is approximately independent of the amplitude---this is a manifestation of the isochronism of nonlinear systems in the small-amplitude limit, and not strict conservation of energy. Newton unified these observations into the laws of motion: in an inverse-square force field, conservation of angular momentum and conservation of energy jointly constrain the orbit to be a closed ellipse.

Once conservation laws have been established, symmetry emerges naturally. If a system possesses a conserved quantity, this means that there exists some transformation under which the behaviour of the system remains unchanged. Translation in time does not change energy, translation in space does not change momentum, and rotation does not change angular momentum. At the beginning of the twentieth century Noether's theorem gave this intuition a profound mathematical form: \textbf{on the premise that the action principle holds and that the symmetry transformations are continuous and smooth}, every continuous symmetry implies a conserved quantity, that is, \textbf{symmetry $\Rightarrow$ conservation}; \textbf{the converse fails}---a conservation law does not imply the corresponding symmetry (see \cref{cor:conservation_symmetry_hierarchy} of the `Operations and Periodicity' chapter). \textbf{The conditions of Noether's theorem cannot be omitted}---for discrete symmetries, for non-Lagrangian systems or for non-smooth systems the above implication generally fails, and the cases have to be treated separately.

The logical order must be stressed: \textbf{it is not that symmetry comes first and conservation laws afterwards; rather, the observation of invariance comes first, conservation laws are then abstracted from it, and symmetry is discovered after that}. Within the framework of Noether's theorem, conservation laws and symmetries do not stand in one-to-one correspondence but in a \textbf{one-directional implication}: \textbf{symmetry $\Rightarrow$ conservation}, whilst the converse fails, that is, \textbf{conservation $\nRightarrow$ symmetry} (\cref{cor:conservation_symmetry_hierarchy}). They point to one and the same fact: the system preserves a certain constant structure throughout change. The manifestation of this structure in phase space is, \textbf{under certain conditions}, a closed orbit or an invariant torus. \textbf{But a conservation law by itself does not guarantee that the orbit is closed}---a system with conservation of energy may still have an orbit that does not close (a hyperbolic orbit, for instance); closure requires additional topological or dynamical conditions.

\section*{Three generating modes of periodicity: a research programme}

Within the framework of this book we \textbf{reduce the generating transformations of periodicity to three basic operations}, and maintain that they constitute a reasonable research programme. The strict form of this claim---namely `any continuous transformation that preserves some intrinsic structure can be decomposed into a composite of these three groups'---is a \textbf{research programme}; the range within which it holds has to be delimited chapter by chapter, and it should not at present be regarded as a proved theorem.

\begin{enumerate}[leftmargin=2em, label={(\arabic*)}]
    \item \textbf{Translation}. On a train travelling at constant speed, each full revolution of a wheel brings some mark on the carriage floor back to the same position. What translation generates is the traditional cyclic period, and its conserved quantity \textbf{within the Lagrangian framework} is momentum.
    \item \textbf{Scaling}. The boundary of the Koch snowflake returns to itself under every scaling transformation---after magnification by a factor of three, the boundary curve has the same structure as before magnification. This is a \textbf{self-similar} structure: the perimeter is infinite whilst the area is finite, and the area is precisely the conserved quantity that makes scaling `close'.
    \item \textbf{Rotation}. The revolution of the Earth about the Sun is, in the approximation that neglects the perturbations of the other planets, a closed orbit in a rotationally symmetric system, and its conserved quantity is angular momentum.
\end{enumerate}

The three operators gradient, curl and divergence are each given a rigorous definition, differentiability preconditions and range of applicability in the `Vector Analysis' chapter; their respective structures have already been listed in the \textbf{Introduction} and are not repeated here.

\textbf{It should be pointed out in particular} that this book does \textbf{not} place these three operators alongside the three basic transformations of translation, scaling and rotation as a set of `correspondences'. The three transformations act on \textbf{space}, whilst the three operators act on \textbf{fields} (scalar fields, vector fields), and the two lie at different mathematical levels; arranging them into a correspondence table is both unnecessary and liable to obscure their respective genuine preconditions. Any statement involving any one of the three is given separately in the `Vector Analysis' chapter, with explicit differentiability preconditions and domain conditions.

\section*{Phase: the mathematics of relation itself}

In the mathematics of periodicity, phase \textbf{can} be understood as an active transformation operator, and not merely as a passive coordinate. The $k$-th phase of a period $T$ corresponds to the complex number $\mathcal{O}_k^{(T)} = e^{i2\pi k/T}$. Multiplication by $\mathcal{O}_0 = 1$ yields the object itself---congruence; multiplication by $\mathcal{O}_{T/2} = -1$ yields the opposite vector---antipodality; multiplication by an arbitrary $\mathcal{O}_k$ defines the relation of rotation through $2\pi k/T$. The multiplicative closure of the phase operators guarantees the closure of the system of relations: a relation between relations is still a relation. This closure is an algebraic property of the finite cyclic group $C_T$, and does not depend on any physical interpretation.

This perspective offers a \textbf{possible} route towards incorporating such traditional relations of the \textit{I Ching} as `combination and transformation' and `clash and overcoming' into the framework of group representation theory; its concrete modern mathematical encoding (the correspondences of the Heavenly Stems, the Earthly Branches and the Eight Trigrams with groups with periodic orbits, together with the epicycle nesting rules involved) is given in the \textbf{Introduction}, where the book formulates these correspondences as algebraic isomorphisms or equivalences of group actions, rather than as historical assertions about ancient documents.

\section*{Leibniz's discovery: periodicity embedded in number itself}

In the course of the mathematisation of periodicity, Leibniz made a distinctive and profound observation. Whilst studying binary arithmetic he discovered that, in the binary expansions of the sequence of natural numbers, the digit in each position displays a strictly periodic behaviour---the $0$-th position cycles with period $2$, the $1$-st position cycles with period $4$, and the $k$-th position cycles with period $2^{k+1}$.

\textbf{A note on historical verification}: Leibniz did study binary arithmetic and was deeply interested in its philosophical implications (in particular its connection with the \textit{I Ching}); the sentence in this book, `he named this phenomenon \textbf{Periodus numerorum}---the period of the natural numbers', belongs to the category of \textbf{historical material awaiting verification}, and would require the support of primary sources or of secondary scholarship before it could be stated as historical fact. Here we retain this formulation, but we mark its epistemological status as `alleged' or `historical material awaiting verification', and do not treat it as a settled historical conclusion.

The \textbf{mathematical content} of this observation is itself clear: periodicity does not exist only in the grand motions above the firmament, but is also embedded in the structure of the binary representation of the number system. In each position the $0$s and $1$s alternate with a fixed rhythm.

Leibniz went further and pointed out that the uniform distribution of binary fractions on the unit circle of the complex plane---taking $n$ binary places divides the circumference into $2^n$ equal points---constitutes an approximation to the continuous circle. This geometrical intuition \textbf{can} be set in correspondence with the generative structure of the hexagram system of the \textit{I Ching}: Taiji ($2^0$) $\to$ Yin-Yang ($2^1$) $\to$ the Four Symbols ($2^2$) $\to$ the Eight Trigrams ($2^3$) $\to$ the Sixty-Four Hexagrams ($2^6$). \textbf{It should be pointed out} that this correspondence is a modern mathematical interpretation offered by this book, and not a historical assertion about the original text of the \textit{I Ching}.

\section*{Epicycles: a unified language for describing periodic phenomena}

In systems possessing conservation laws there exist closed orbits \textbf{under certain conditions}. A planet with conservation of angular momentum has an orbit that closes in an inverse-square force field; a pendulum with conservation of energy has a swing that is a closed orbit in phase space. A closed orbit means that the system traces a loop in phase space---an epicycle.

The epicycle is not the exclusive property of astronomy. \textbf{Many} things with a periodic structure can all be described by means of epicycles. The propagation of a beam of monochromatic light can be described by the sinusoidal oscillation of the electromagnetic field in space; the acoustic vibration of a word can be approximately described by the periodic fluctuation of air pressure; the prosody of a sentence can be approximately described by the periodic alternation of syllable length and stress; the oscillation of an electrical signal in a wire can be approximately described by the back-and-forth motion of charge about an equilibrium position; and the crystal oscillator in a chip is a substantially fixed epicycle that provides the whole system with its time reference.

Mathematically speaking, \textbf{under suitable function spaces and convergence conditions}, a periodic function can be represented as a Fourier series:
\[
f(x) = \sum_{k \in \mathbb{Z}} c_k e^{i k \omega_0 x}.
\]
\textbf{Two cases have to be distinguished}: (i) the \textbf{finite-harmonic model}: $f(x) = \sum_{k=1}^{N} A_k e^{i(\omega_k x + \phi_k)}$, suitable for describing the superposition of finitely many epicycles, and a form in common use in numerical implementation and engineering modelling; (ii) the \textbf{general Fourier series}: $f(x) = \sum_{k \in \mathbb{Z}} c_k e^{i k \omega_0 x}$, the complete expansion of a periodic function (under suitable convergence conditions), whose coefficients are complex numbers and carry amplitude and phase information. \textbf{One cannot say that `a periodic function = a superposition of finitely many epicycles'}---a general periodic function requires infinitely many Fourier modes (a square wave, for example, requires infinitely many odd harmonics). But when attention is concentrated on finitely many dominant components, the finite-harmonic model is a useful approximation. When concrete numerical implementation is involved, this book usually adopts the finite-harmonic model, whilst when general theory is involved it uses the Fourier series.

Under appropriate conditions, periodic phenomena can be understood as superpositions of several epicycles. The large circle of the outer ring carries the overall rhythm, the small circles of the inner rings contribute the texture of detail, and the phase determines the precise position of each circle at a given moment.

\section*{Periodic convolution: a unified computation for cross-period interaction}

One of the core research programmes proposed in this book is \textbf{periodic convolution}: bringing amplitude (scaling) and phase (rotation) into the operations on series as well, and thereby retaining the interaction between every pair of frequencies; the second line of formulae on the cover is its preliminary form. Its two-layer construction (the interaction tensor and the composition map), the periodic convolution formed by composing the two, and the various conditions required for it to hold rigorously, are given together in the \textbf{Introduction} and are not repeated here.

\section*{Approximation, intercalation and the boundedness of incommensurable periods}

Among the many technical difficulties confronting the mathematics of periodicity, one class is unlike the others: it comes neither from computation nor from topology, but from \textbf{number theory}. When the continuous period $T_f$ and the sampling interval $\Delta x$ are incommensurable, that is, when the ideal period ratio
\[
\rho = \frac{T_f}{\Delta x} \notin \mathbb{Q}
\]
then the sampling sequence within any finite window \textbf{cannot} have a strictly integral period. This is not an observational error, nor a defect of the instrument, but the intrinsic cost of incommensurability. The manifestation of this cost in the frequency domain is \textbf{spectral leakage}: the number of periods within the window is not an integer, and the energy of the signal necessarily spreads from the main lobe into the sidelobes.

\subsection*{The mathematical definition of intercalation}

One of the core methodologies proposed in this book is \textbf{intercalation}---approximating the real number $\rho$ by a sequence of integers $\{N_m\}$ so that the \textbf{cumulative sample count}
\[
S_M = \sum_{m=0}^{M-1} N_m
\]
approximates $M\rho$ with \textbf{bounded error}. The core construction of intercalation consists of three basic intercalation sequences:
\begin{enumerate}[leftmargin=2em, label={(\arabic*)}]
    \item \textbf{Floor intercalation}: $N_m^{\mathrm{floor}} = \lfloor (m+1)\rho \rfloor - \lfloor m\rho \rfloor$, with error $E_M^{\mathrm{floor}} \in (-1,0]$;
    \item \textbf{Ceiling intercalation}: $N_m^{\mathrm{ceil}} = \lceil (m+1)\rho \rceil - \lceil m\rho \rceil$, with error $E_M^{\mathrm{ceil}} \in [0,1)$;
    \item \textbf{Balanced intercalation}: $N_m^{\mathrm{round}} = \operatorname{round}((m+1)\rho) - \operatorname{round}(m\rho)$, with error $|E_M^{\mathrm{round}}| \le 1/2$.
\end{enumerate}
What the three sequences have in common is this: $N_m \in \{\lfloor\rho\rfloor, \lceil\rho\rceil\}$---at each step the choice is only between the lower integer period and the upper integer period---yet the cumulative error is confined within a bounded band, either $[-1/2, 1/2]$, $(-1,0]$ or $[0,1)$. This boundedness is where the whole value of intercalation lies.

\subsection*{Foundations of approximation theory: the role of continued fractions}

The \textbf{optimal} choice of intercalation sequence is not arbitrary; it is determined by the continued-fraction expansion of $\rho$. Let $\{p_j/q_j\}$ be the continued-fraction convergents of $\rho$. Then:
\begin{enumerate}[leftmargin=2em, label={(\roman*)}]
    \item \textbf{Optimality}: $p_j/q_j$ is the closest approximation to $\rho$ among all rational numbers whose denominators do not exceed $q_{j+1}-1$;
    \item \textbf{Error bound}: $|p_j - q_j\rho| < 1/q_{j+1}$;
    \item \textbf{Alternating approximation}: the even convergents approach $\rho$ from below and the odd ones from above, forming a symmetric squeeze.
\end{enumerate}
On the intercalation window $T = p_j$, $M = q_j$, the cumulative error tends to zero. This raises intercalation from an `error-control trick' to a \textbf{method with a rigorous optimality criterion}---continued-fraction convergents are precisely the optimal intercalation windows. This optimality derives from a classical fact in number theory: continued-fraction convergents are the best rational approximations to a real number; intercalation translates this from number theory into the choice of sampling window.

\subsection*{Bounding spectral leakage}

The effect of intercalation on spectral leakage can be stated precisely. Let $F[k]$ be the DFT on the intercalation window and $c_m$ the continuous Fourier coefficients. Then
\[
\left| F[k] - c_{k/M} \right| \le C \cdot \frac{|E_M|}{T} \cdot \sum_{m\in\mathbb{Z}} |c_m|,
\]
where $C$ is a constant independent of $T$. This inequality gives an \textbf{explicit upper bound} for the leakage---intercalation reduces leakage from `uncontrollable' to `bounded'. When $\rho \in \mathbb{Q}$, intercalation degenerates into a constant sequence and the leakage is strictly zero; when $\rho \notin \mathbb{Q}$, intercalation chooses windows by means of continued fractions, so that this upper bound tends to zero along the sequence of convergent windows.

\textbf{It must be stressed} that intercalation does \textbf{not eliminate} leakage; it reduces the amplitude of the leakage from $O(1)$ to $O(|E_M|/T)$. This distinction is essential, and it reflects the intrinsic cost of incommensurability---no finite window can align $\rho$ exactly with an integer. The value of intercalation lies in this: it provides a \textbf{quantitative characterisation} of this cost and a \textbf{systematic method of minimising} it.

\subsection*{The complementarity of intercalation and window functions}

Intercalation and window functions are two independent routes for suppressing spectral leakage, and they act on different stages of the leakage:
\begin{itemize}[leftmargin=2em]
    \item \textbf{Intercalation} acts on the \textbf{window length}: choosing $T = p_j$ makes the window contain approximately an integral number of continuous periods, reducing leakage at its source;
    \item \textbf{Window functions} act on the \textbf{window boundary}: smoothing windows (Hann, Hamming, Blackman and the like) suppress the sidelobes, at the cost of a broadening of the main lobe.
\end{itemize}
The two can be used in combination: first use intercalation to choose the optimal window, then apply a window function, and one obtains a narrow main lobe and low sidelobes simultaneously. This is a combined effect that no single method can achieve.

\subsection*{The place of intercalation in this book}

The method of intercalation performs three functions in this book:
\begin{enumerate}[leftmargin=2em, label={(\arabic*)}]
    \item in the `Discrete Sampling and Intercalation of Periodic Signals' chapter it is established as a bridge between commensurability and incommensurability, yielding the intercalation error boundedness theorem and the continued-fraction optimal intercalation theorem;
    \item in the `Harmonic Analysis' chapter it is used to derive the DFT leakage estimate for intercalation windows, giving an explicit upper bound for the amplitude of the leakage;
    \item in the `Periodic Category' chapter it is made functorial---the intercalation functor $\mathcal{R}$ repairs the irreversible gap from the continuous category to the discrete category, so that the adjoint pair formed by the continuous limit and the discretisation functor is preserved to arbitrary accuracy in the incommensurable case.
\end{enumerate}
Furthermore, intercalation, together with the window-function methods of the `Time Series Analysis' chapter, calendrical intercalation in the `Periodic Structure of Astronomical Systems' chapter, intercalation by absence of mid-terms in the `Periodic Structure of Calendar Systems' chapter, state periods in the `Periodic Structure of Cryptographic Systems' chapter, multi-agent synchronisation in the `Periodic Structure of Agents' chapter, and multitaper spectral estimation in the `Periodic Structure of Data Analysis' chapter, forms a unified perspective that runs through the whole book.

\textbf{It must be stressed once again} that intercalation is at present established throughout the book as a methodology supported by rigorous theorems, but that it does \textbf{not solve} all the difficulties of incommensurability. In cases such as multi-frequency signals, time-varying periods and unknown period ratios, the role of intercalation is limited, and it has to be used in combination with other methods (multitaper spectral estimation, Bayesian spectral analysis, joint continued-fraction approximation). This book will \textbf{clearly mark} the range of applicability and the limitations of the intercalation method.

\section*{From disciplinary fragmentation to a unified perspective}

In the existing system of disciplines, researchers in different fields all deal with periodic phenomena, yet they use mutually disconnected languages. Astronomers use perturbation theory and orbital resonance to describe the quasi-periodic motions of celestial bodies; condensed-matter physicists use space groups and Brillouin zones to describe the translational symmetry of crystals; biologists use limit cycles and phase response curves to describe the genetic networks of circadian rhythms; fractal geometers use self-similarity and scaling exponents to describe invariant structures across scales; economists use spectral analysis and autoregressive models to detect the characteristic frequencies of business cycles; and computer scientists use cyclic redundancy checks and fast algorithms to handle the periodic components of signals.

Every one of these researchers is dealing with a problem connected with periodicity, yet each uses a mutually disconnected system of terminology. What a condensed-matter physicist calls a `reciprocal lattice vector', what a signal-processing engineer calls the `Nyquist frequency', and what a harmonic analyst calls `Pontryagin duality'---these point, \textbf{within the harmonic analysis of finite cyclic groups or of discrete translation groups}, to one and the same mathematical fact: the frequency-domain representation of a periodic function. \textbf{It should be pointed out} that this `sameness' holds at an appropriate level of abstraction, and that the concrete physical meaning and the conditions of applicability may differ from one discipline to another.

A similar fragmentation exists in the domain of integral transforms. The Fourier transform, the Laplace transform, the wavelet transform and the Mellin transform are usually taught separately, as mutually independent tools. Yet when examined from the standpoint of group representation theory, \textbf{if appropriate group actions, representations and function spaces are chosen}, the different integral transforms can be given a unified description within the framework of group representations and harmonic analysis: the Fourier transform is associated with the translation group, the Mellin transform with the scaling group, and the wavelet transform with the affine group. \textbf{These statements need to be verified case by case}, and their precise form depends on the chosen function space, group action and measure.

In sampling theory this fragmentation likewise exists. The `sampling theorem' in signal-processing textbooks, `Pontryagin duality' in harmonic-analysis textbooks, and `intercalation' in calendrical-astronomy textbooks---the three are taught separately, yet they point to three faces of one and the same mathematical structure: how, in the incommensurable case, to approximate the continuous by the discrete. The theory of intercalation presented in this book is precisely an attempt to bring the three within a single framework.

\section*{The cost of disciplinary barriers}

The fragmentation of terminology and method described above brings at least threefold costs:

\begin{enumerate}[leftmargin=2em, label={(\arabic*)}]
    \item \textbf{Repeated discovery}. A mathematical theorem already known in one discipline is independently rediscovered in another, resulting in the duplicated investment of scholarly resources.
    \item \textbf{Deep analogies are obscured}. Between the Kepler ellipse and the Koch snowflake, between crystal lattices and genetic oscillations, between the Fourier transform and the Mellin transform---these identical mathematical structures, which hold \textbf{at an appropriate level of abstraction}, are difficult to recognise because of differences in disciplinary language.
    \item \textbf{Lack of diagnostic tools}. Faced with a new and complex system, a researcher may be able to rely only on descriptions at the level of phenomena, and finds it difficult to carry out a diagnosis at the level of structure.
\end{enumerate}

\section*{The possibility of a unified framework}

Although the disciplines differ on the surface, in every case involving a periodic structure it is \textbf{usually} possible to isolate three common mathematical ingredients:

\begin{itemize}[leftmargin=2em]
    \item a \textbf{symmetry group} $G$, acting on the state space $M$ of the system;
    \item a \textbf{conserved quantity} $Q$, invariant under the action of $G$;
    \item a \textbf{closed orbit} $\mathcal{O} \subset M$, generated by the joint constraints of the conservation laws and the dynamics.
\end{itemize}

This means that a \textbf{unified mathematical framework of periodicity}, whose language is group action, whose criterion is a conserved quantity, whose geometrical image is a closed orbit and whose tool of decomposition is the epicycle series, is theoretically possible. This book is precisely one exploration of that possibility. \textbf{Whether it succeeds depends on whether it can produce non-trivial new mathematical results in the chapters that follow}, and not merely on renaming knowledge that already exists.

In the chapters that follow, this book will gradually develop the complete theory of the intercalation method: in the `Discrete Sampling and Intercalation of Periodic Signals' chapter as a tool of commensuration, in the `Harmonic Analysis' chapter as a tool for bounding leakage, and in the `Periodic Category' chapter as a tool of functorial correction. Intercalation turns `incommensurability', the most stubborn number-theoretic difficulty in periodic analysis, from a passively accepted defect into an actively exploited structure---this is an independent contribution of this book to the methodology of periodic analysis.

\section*{The standpoint and the limits of this book}

The mathematics of periodicity does not claim a priori that all the systems it examines `possess a periodic structure'. In each application chapter it begins by asking three questions, by means of three criteria (the full statement of the criteria is given in the \textbf{Introduction}). For systems that satisfy the conditions, the book gives a rigorous description of the periodic structure; for systems that satisfy them only in part, the book measures precisely the degree of their deviation from a closed orbit; for systems that do not satisfy them, the book states clearly that their `periodicity' is merely a metaphor at the level of phenomena. This diagnostic function is the key respect in which the mathematics of periodicity differs from general interdisciplinary analogy.

Nor does the mathematics of periodicity attempt to become a theory of everything. It does not explain why a particular conserved quantity takes a particular value, does not derive the evolution equation of a particular system, and does not predict the future course of the economy. What it provides is a set of diagnostic tools---a unified grammar for rigorously describing, comparing and analysing the periodic structures that arise in different systems. Its success lies not in proving that everything is periodic, but in being able to state precisely: this system is periodic, because its orbit closes under the action of some group; that system appears to be periodic, but in fact lacks the guarantee of a conservation law; this integral transform and that integral transform are equivalent representations of harmonic analysis on the same group under different parametrisations. \textbf{The correctness of these judgements has to be supported, chapter by chapter, by concrete mathematical propositions.}

\section*{The journey of this book}

The journey of this book sets out from the simplest translation period, gradually builds up the mathematical theory of rotation periods and scaling periods, passes through group representation theory, harmonic analysis and periodic manifolds, and finally arrives at the abstract unification of the periodic category. Once this framework has been established, we shall take this set of tools with us and travel in turn through astronomy, calendrics, geography, physics, biology, psychology, economics, political science, computer science and integral transforms---in each field revealing, \textbf{as far as possible} in the form of theorems or rigorous propositions, how the mathematical structure of periodicity illuminates our understanding of nature and society; and at the same time, \textbf{for those parts where a rigorous proposition cannot yet be given, the book will clearly mark them as conjectures or analogies}.

That nested picture of epicycles on the cover is precisely the starting point of this journey. The points on the three concentric circles are the outermost celestial sphere; the nested small circles extending from the red point at the centre are the trajectories of the planets; and the two lines of formulae on the cover are the mathematical language of all of it. When we close this book, we hope that what the reader sees will no longer be merely periodic functions, but a living, vibrant universe woven together out of amplitude, frequency, phase and period---\textbf{provided, of course, that we always remember that what mathematics describes is the model, and not the universe itself}.

\section*{Intended readership and prerequisites}

This book addresses four kinds of reader: mathematicians who wish to treat the periodic structure of a system as an algebraic and geometrical object; graduate students and researchers in science and engineering who deal with periodic phenomena daily and wish to acquire a unified vocabulary; engineers and data analysts who face incommensurable sampling ratios and spectral leakage in their practical work; and readers with a general mathematical literacy who are concerned with the relation between modern mathematics and traditional symbolic systems (such as the \textit{I Ching}).

The prerequisites required by the mathematical chapters are: linear algebra (matrices, eigenvalues, transposition, the orthogonal group); mathematical analysis (limits, continuity, differentiability, uniform convergence); group representation theory (groups, homomorphisms, orbits, group actions, representations); and elementary complex analysis (the complex plane, the exponential function, holomorphic functions). The two chapters that require topology and differential geometry---`Manifold and Fibre Bundle' and `Periodic Manifold'---each declare their own prerequisites at the beginning of the chapter; readers who lack this preparation may skip these two chapters on a first reading without affecting the logical thread of the remaining parts. Apart from the first eleven chapters, the application chapters require no further mathematics, and each chapter restates the concepts it uses; the `Probability and Distribution' and `Time Series Analysis' chapters additionally require the foundations of an introductory course in probability theory and mathematical statistics.

\textbf{Reading paths}: readers whose main interest is the mathematical core may read Chapters 1 to 24, together with Chapter 46; readers whose main interest is one particular application may read Chapter 1, then the relevant chapter in Part II, and then the application chapter that concerns them; readers who need only computational tools may begin with the four chapters `Discrete Sampling and Intercalation of Periodic Signals', `Harmonic Analysis', `Epicycle Activation Function and Function Approximation' and `Modelling and Analysis: From Data to Epicycles'.

\section*{Conventions of style and labelling}

The mathematical chapters of this book follow this order: \textbf{Definition} $\to$ \textbf{Proposition} $\to$ \textbf{Lemma} $\to$ \textbf{Theorem} $\to$ \textbf{Proof} $\to$ \textbf{Corollary} $\to$ \textbf{Property} $\to$ \textbf{Example} $\to$ \textbf{Remark}. Cross-references are always given by name, and the referencing mechanism of the whole book uses the standard facilities provided by the \texttt{cleveref} package of \LaTeX{}.

\textbf{The proof status of every assertion is marked}. Those for which a proof is given in this book are presented in the form of a definition, proposition, lemma, theorem, corollary or property, and are accompanied by a proof. Those for which no proof is given in this book are marked with one of the following three labels: \textbf{Conjecture}, meaning that the statement is precise but that no proof has yet been obtained; \textbf{Analogy}, meaning that the correspondence between two structures is suggestive, but is not a bijection; \textbf{Research programme}, meaning a framework for organising a class of problems, rather than the assertion of a theorem. \textbf{These three labels should be read as warnings, not as ornament}: any claim bearing one of them has no proof available, and this book does not use it as a premise in any derivation.

Notation: $\mathbb{N}, \mathbb{Z}, \mathbb{Q}, \mathbb{R}, \mathbb{C}, \mathbb{H}$ denote respectively the natural numbers, the integers, the rational numbers, the real numbers, the complex numbers and the quaternions. The three centres of symmetry are written $c=0$ (additive reflection), $c=1$ (multiplicative scaling reflection) and $c=i$ (multiplicative inversion). The period is written $T$, the corresponding angular frequency $\omega = 2\pi/T$, and the phase operator $\mathcal{O}_k^{(T)} = e^{i2\pi k/T}$. Periodic convolution is written $\circledast$; when the order $d$ or the window kernel has to be displayed, it is written $\mathcal{K}^{(d)}\{f,g\}$.

\textbf{Concerning historical sources}: the historical and cultural material in the preface and in the `Periodic Structure of the I Ching System' chapter falls into two classes. The parts concerning traditional calendrical reckoning, the \textit{I Ching} and the history of astronomy are all cited as \textbf{traditional accounts or alleged historical sources}; they are not premises of any mathematical derivation in this book, and their specific numerical values should be regarded as symbolic settings. The parts concerning modern mathematics are given together with their premises. Wherever primary sources are required but cannot be supplied by this book, the text marks the passage as `historical material awaiting verification'.

\section*{Acknowledgements}

This book originated in an attempt to write down, in one and the same notation, the periodic structures that appear in fields as far apart as celestial mechanics, crystallography, signal processing and calendrical reckoning. In the course of doing so it has incurred many debts. It owes a debt first of all to the mathematicians whose theorems are cited in this book, and to the authors of the textbooks and monographs listed in the bibliography---the standard results of group representation theory, harmonic analysis, number theory and differential geometry that appear in this book are all theirs, and whatever is new in this book has been obtained by standing on those results.

It owes a particular debt to the classical Chinese literature on mathematics and astronomy---including the \textit{Huainanzi}, the \textit{Huangji Jingshi}, the twenty-four solar terms and the sexagenary cycle: this literature provides the earliest instances of the common-multiple reasoning that this book calls `intercalation'. \textbf{It must be stressed once again} that the modern mathematical reading offered in this book is the author's own interpretation, and not a historical assertion about these texts.

Thanks are due to the maintainers of the open-source tools without which this book could not have been completed: \TeX{} Live and the \LaTeX{} packages used for typesetting, and the developers of the free typefaces used for the Chinese fonts. The errors that remain are the author's own.

Corrections, counterexamples and refutations are \textbf{especially} welcome. Since the book is offered as a research programme rather than as a finished edifice, the most useful response to it is to point out where it fails---and as precisely as possible.

\vspace{1em}
\begin{flushright}
Independent Research in Periodicity Mathematics\\
September 2026
\end{flushright}

